\SetChapterInfo{subtitle={Schwartz 对偶与微分的乘子化},keywords={温和分布; Fourier 自同构; Sobolev 空间},prereq={测试函数拓扑; 分布运算与卷积; 经典 Fourier 反演},goals={统一常数并使用对偶定义 Fourier 变换}}
\chapter{温和分布与傅里叶变换}\label{ch:fourier}
\ChapterSource{对应施瓦尔茨 \cite{schwartz} 第 VII 章 §§2--6，特别是原书印刷第 169、172 页。全章使用式 \eqref{eq:fourier-convention}，与原书频率归一化不同。}

\section{温和分布的增长控制}
\begin{definition}[温和分布]
Schwartz 空间 $\Sch(\R^n)$ 的连续线性对偶记为 $\Temp(\R^n)$。其元素称为温和分布，也称缓增分布。
\end{definition}
由有限半范数估计，$T\in\Temp$ 当且仅当存在 $N,m,C$ 使
\begin{equation}\label{eq:tempered-bound}
 |T(\varphi)|\le Cq_{N,m}(\varphi)\qquad(\varphi\in\Sch).
\end{equation}
$\Test\to\Sch$ 连续，故限制得到 $\Temp\to\Dist$；由于 $\Test$ 在 $\Sch$ 中稠密，这一映射单射。也就是说，温和分布是普通分布中的特殊一类，它们能连续地作用于非紧支撑但快速衰减的探针。

\begin{proposition}[基本例子]
以下对象都是温和分布：$L^p(\R^n)$ 函数（$1\le p\le\infty$）；具有至多多项式增长的局部可积函数；紧支集分布；$\delta_a$ 的任意阶导数。
\end{proposition}
\begin{proof}
对 $L^p$ 函数用 Hölder 不等式，再以足够高权重控制测试函数的 $L^{p'}$ 范数。若 $|f(x)|\le C(1+|x|)^M$ 几乎处处，取 $N>M+n$ 得
\[
 |\int f\varphi|\le Cq_{N,0}(\varphi)
 \int(1+|x|)^{M-N}\,\dd x<\infty.
\]
紧支集分布的有限阶估计由某个 $q_{0,m}$ 控制，故也温和。
\end{proof}
更一般，若 $\int|f(x)|(1+|x|)^{-N}\,\dd x<\infty$，则 $T_f$ 温和。对带符号或复值的振荡函数，这只是充分条件，不能误当成必要条件。

\begin{example}[普通分布并非都温和]
函数 $e^x$ 在 $\R$ 上局部可积，给出分布但不是温和分布。若它满足式 \eqref{eq:tempered-bound}，取非负非零 $\eta\in\Test$ 并令 $\eta_j(x)=\eta(x-j)$，则配对为 $e^j\int e^y\eta(y)\,\dd y$，而 $q_{N,m}(\eta_j)$ 至多按 $j^N$ 增长，矛盾。
\end{example}
求导和多项式乘法保持 $\Temp$。更一般，若光滑 $a$ 的每个导数均有某种多项式增长，则 $\varphi\mapsto a\varphi$ 连续作用于 $\Sch$，其转置给出温和分布乘法。

\section{Schwartz 空间上的傅里叶变换}
\begin{theorem}[Fourier 自同构]\label{thm:schwartz-fourier}
按式 \eqref{eq:fourier-convention} 定义的 $\FT$ 是 $\Sch(\R^n)$ 的连续线性自同构，其逆为
\[
 \FT^{-1}\psi(x)=(2\pi)^{-n}\int e^{ix\cdot\xi}\psi(\xi)\,\dd\xi.
\]
\end{theorem}
\begin{proof}
微分可移入积分，分部积分无边界项，故
\begin{equation}\label{eq:fourier-identities}
 \widehat{\partial_j\varphi}=i\xi_j\widehat\varphi,
 \qquad\widehat{x_j\varphi}=i\partial_{\xi_j}\widehat\varphi.
\end{equation}
任意 $\xi^\beta\partial_\xi^\alpha\widehat\varphi$ 因此可表示成 $x^\alpha\varphi$ 的有限阶导数的 Fourier 变换，差一个常数因子。其上确界由对应函数的 $L^1$ 范数控制，而该范数由有限个 $q_{N,m}(\varphi)$ 控制。故 $\FT$ 连续映入 $\Sch$。逆号积分同理连续映入 $\Sch$。使用经典 Schwartz 函数 Fourier 反演公式，两个映射互为逆。
\end{proof}
经典反演也可以用高斯截断验证：在逆积分中加入 $e^{-\varepsilon|\xi|^2}$，得到与高斯核的卷积；高斯核构成近似恒等，令 $\varepsilon\downarrow0$ 得到原函数。这里所需 Fubini 与控制收敛均由 Schwartz 衰减保证。

\section{温和分布的傅里叶变换}
\begin{definition}[对偶定义]
对 $T\in\Temp$，定义
\begin{equation}\label{eq:fourier-distribution}
 \Pair{\widehat T}{\varphi}=\Pair{T}{\widehat\varphi},\qquad\varphi\in\Sch.
\end{equation}
\end{definition}
注意右侧是探针的正向变换。这是因为核 $e^{-ix\cdot\xi}$ 在两个变量中对称，配对采用双线性约定。若 $T=T_f$ 且 $f\in L^1$，Fubini 定理给出该定义与普通 $\widehat f$ 一致。连续自同构的转置也是连续自同构，因此 $\widehat T\in\Temp$ 且 Fourier 变换可逆。
\begin{keyformula}[常数与基本对象]
\begin{align*}
 \widehat{\delta_0}&=1,&\widehat1&=(2\pi)^n\delta_0,\\
 \widehat{\delta_a}&=e^{-ia\cdot\xi},&
 \widehat{\partial^\alpha T}&=(i\xi)^\alpha\widehat T,\\
 \widehat{x_jT}&=i\partial_{\xi_j}\widehat T,&
 \widehat{\tau_hT}&=e^{-ih\cdot\xi}\widehat T.
\end{align*}
\end{keyformula}
例如 $\widehat{\delta_0}(\varphi)=\widehat\varphi(0)=\int\varphi$。又由反演 $\int\widehat\varphi=(2\pi)^n\varphi(0)$，得到常数的变换。对于导数，
\[
 \widehat{\partial_jT}(\varphi)
 =-T(\partial_{x_j}\widehat\varphi)
 =T(\widehat{i\xi_j\varphi})
 =\widehat T(i\xi_j\varphi).
\]
这直接验证 $i$ 与负号，无须靠记忆。

\begin{example}[高斯]
对 $t>0$，令
\[
 G_t(x)=(4\pi t)^{-n/2}e^{-|x|^2/(4t)}.
\]
一维高斯积分与张量积计算给出 $\widehat G_t(\xi)=e^{-t|\xi|^2}$。由于 $\int G_t=1$，$G_t\to\delta_0$ 于分布意义；变换侧则 $e^{-t|\xi|^2}\to1$ 于温和分布意义。
\end{example}

\Needspace{10\baselineskip}
\section{卷积定理的合法范围}
对 Schwartz 函数 $f,g$，Fubini 直接给出 $\widehat{f*g}=\widehat f\widehat g$。若 $S\in\CompactDist$、$T\in\Temp$，则 $S*T\in\Temp$，并且
\begin{equation}\label{eq:convolution-fourier}
 \widehat{S*T}=\widehat S\,\widehat T.
\end{equation}
紧支集分布的 Fourier 变换是光滑函数
\[
 \widehat S(\xi)=\Pair{S_x}{e^{-ix\cdot\xi}},
\]
这里以紧支集切断解释非紧支撑探针。对 $\xi$ 求导产生 $(-ix)^\alpha$；有限阶估计表明每个导数至多多项式增长。因此它是合法 Schwartz 乘子，右端有意义。卷积定理可由对探针的配对与第 \ref{ch:convolution} 章的切断定义验证。

任意两个温和分布的卷积未必存在，它们的 Fourier 变换之积也未必有定义。例如 $1*1$ 无定义，相应右侧形式上含 $\delta_0^2$。在允许的函数类中，乘积定理为 $\widehat{fg}=(2\pi)^{-n}\widehat f*\widehat g$；其中常数因子来自反演。

\section{Sobolev 空间：分布的 Hilbert 尺度}
\begin{definition}[Bessel--Sobolev 空间]
对 $s\in\R$，$H^s(\R^n)$ 由那些 Fourier 变换可用可测函数表示、且满足下列有限范数条件的 $u\in\Temp$ 组成：
\[
 \|u\|_{H^s}^2=(2\pi)^{-n}\int_{\R^n}
 (1+|\xi|^2)^s|\widehat u(\xi)|^2\,\dd\xi<\infty.
\]
\end{definition}
Plancherel 定理给出 $H^0=L^2$。对非负整数 $m$，权重 $(1+|\xi|^2)^m$ 与 $\sum_{|\alpha|\le m}|\xi^\alpha|^2$ 等价，故 $H^m$ 等价于所有不超过 $m$ 阶弱导数属于 $L^2$ 的空间。并有
\[
 \partial_j:H^s\to H^{s-1}\quad\text{连续},\qquad
 s_1\ge s_2\Rightarrow H^{s_1}\subset H^{s_2}.
\]
在本稿的双线性分布配对下，$H^{-s}$ 与 $H^s$ 的连续线性对偶的对应明确为
\[
 \langle u,v\rangle=(2\pi)^{-n}\int_{\R^n}
 \widehat u(\xi)\widehat v(-\xi)\,\dd\xi,
 \qquad |\langle u,v\rangle|\le\|u\|_{H^{-s}}\|v\|_{H^s}.
\]
先在 Schwartz 函数上由反演得到公式，再用加权 Cauchy--Schwarz 延拓。反向将 $v\mapsto\langle\xi\rangle^s\widehat v(\xi)$ 作为到 $L^2$ 的同构，对任意连续泛函用 Riesz 表示，再反射频率并共轭表示函数，即得到唯一 $u\in H^{-s}$，且泛函范数等于 $\|u\|_{H^{-s}}$。因此这里的识别是线性的；若改写成 Hilbert 内积形式，必须同时调整共轭位置。
\begin{example}[点质量的正则性门槛]
因 $\widehat{\delta_0}=1$，有
\[
 \delta_0\in H^s(\R^n)\quad\Longleftrightarrow\quad s<-\frac n2.
\]
积分在原点附近有限，在无穷远需 $\int_1^\infty r^{2s+n-1}\,\dd r<\infty$。类似地，$\partial^\alpha\delta_0\in H^s$ 当且仅当 $s<-n/2-|\alpha|$，可在球面上一块多项式不为零的区域给出必要性。
\end{example}

\section{从 Sobolev 范数到普通函数}\label{sec:topic-sobolev-embedding}
\begin{proposition}[全空间的连续嵌入]\label{prop:topic-sobolev-embedding}
若 $k\in\N_0$ 且 $s>k+n/2$，则 $H^s(\R^n)$ 中每个元素有唯一的 $C^k$ 代表元，各阶 $|\alpha|\le k$ 的导数有界且在无穷远趋于零，并满足
\[
 \max_{|\alpha|\le k}\|\partial^\alpha u\|_\infty
 \le C_{s,k,n}\|u\|_{H^s}.
\]
\end{proposition}
\begin{proof}
加权 Cauchy--Schwarz 不等式给出
\[
 \int |\xi|^{|\alpha|}|\widehat u(\xi)|\,\dd\xi
 \le \left(\int(1+|\xi|^2)^s|\widehat u|^2\,\dd\xi\right)^{1/2}
 \left(\int |\xi|^{2|\alpha|}(1+|\xi|^2)^{-s}\,\dd\xi\right)^{1/2}.
\]
最后的积分由 $s>|\alpha|+n/2$ 保证有限。因此 $u$ 及这些导数均由绝对收敛的逆 Fourier 积分表示；控制收敛保证连续性与逐阶微分，Riemann--Lebesgue 引理保证无穷远消失。该函数代表原分布，连续代表元的唯一性由几乎处处相等推出。
\end{proof}
这把分布意义的正则性变成点值和导数估计，与主线第六章 §7 的 Sobolev 引理衔接。严格不等号是本证明中积分收敛的条件，不能直接改为临界等号。区域上的嵌入还需要延拓或局部化；边界迹是另一类结论，见书后补充说明第 \pageref{sec:supp-trace} 页。

\Needspace{12\baselineskip}
\section{习题}
\begin{exercise}[高阶导数的有界性]
设 $u\in H^s(\R^n)$ 且 $s>1+n/2$。证明 $u$ 的连续代表元为全局 Lipschitz 函数，并估计其 Lipschitz 常数。
\end{exercise}
\begin{solution}
上面的嵌入给出 $u\in C^1$ 且 $\|\nabla u\|_\infty\le C\|u\|_{H^s}$。沿 $x$ 到 $y$ 的线段积分，有 $|u(x)-u(y)|\le\|\nabla u\|_\infty|x-y|$。
\end{solution}
\begin{exercise}[二次变换]
证明 $\FT^2T=(2\pi)^n\Reflected T$，其中 $\Reflected T(\varphi)=T(\varphi(-\cdot))$。
\end{exercise}
\begin{solution}
对探针用经典反演得到 $\FT^2\varphi(x)=(2\pi)^n\varphi(-x)$，故 $\FT^2T(\varphi)=T(\FT^2\varphi)=(2\pi)^n\Reflected T(\varphi)$。
\end{solution}
\begin{exercise}[一阶基本解的变换]
令 $E(x)=H(x)e^{-x}$。验证 $(\partial_x+1)E=\delta_0$，并求 $\widehat E$。
\end{exercise}
\begin{solution}
由光滑乘积公式 $E'=\delta_0-E$。又 $E\in L^1$，故 $\widehat E(\xi)=\int_0^\infty e^{-(1+i\xi)x}\,\dd x=(1+i\xi)^{-1}$。变换侧 $(1+i\xi)\widehat E=1$ 与原方程一致。
\end{solution}
